\subchapter{Software Bill of Materials (SBoM)}{Objective:
	Generate and analyze the image's SBoM.
}

\section{Prerequisites}

We are going to use \code{sbom-cve-check} to analyze our SBoM.

\begin{bashinput}
(.venv) $ pip install sbom-cve-check[extra]
\end{bashinput}

\section{The default SPDX2.2 SBoM}

Yocto is able to generate SPDX SBoMs, but the version we are currently using (Scarthgap)
is using the 2.2 standard by default.

This SBoMs are generated during the build of each recipe.
You can find them under \code{build/tmp-glibc/deploy/spdx/2.2}

\begin{bashinput}
$ tree -d build/tmp-glibc/deploy/spdx/2.2/freiheit93/
|-- packages
|-- recipes
`-- runtime
\end{bashinput}

One downside of SPDX2.2 is, that the SBoM is fragmented into a large number of files. This makes it less obvious
to quickly get an overview:

\begin{bashinput}
$ find build/tmp-glibc/deploy/spdx/2.2/freiheit93/  -type f | wc -l
1952
\end{bashinput}

\section{Generating an SPDX3 SBoM}

To be able to generate an SPDX3 SBoM, we need to add the following lines in 
\code{board_setup/meta-security-training/conf/distro/security-training.conf}:

\begin{bashinput}
INHERIT:remove = "create-spdx"
INHERIT:append = " create-spdx-3.0"
\end{bashinput}

(notice the space on the second line)

In order to generate the full SBoM, Yocto needs to rebuild the image (but it will use the cache).

\begin{bashinput}
$ bitbake security-training-image
\end{bashinput}

You should now see a \code{build/tmp-glibc/deploy/spdx/3.0.1} directory.
The "main" SBoM, the one for the image, is under
\code{build/tmp-glibc/deploy/spdx/3.0.1/freiheit93/image/security-training-image-freiheit93-image.spdx.json}

This SBoM is rather short, so it can be reviewed manually, but it is high-level.
However, the SPDX3 directory, \code{build/tmp-glibc/deploy/spdx/3.0.1} includes an SBoM for the rootFS.
That one is not so short:

\begin{bashinput}
$ jq . build/tmp-glibc/deploy/spdx/3.0.1/freiheit93/rootfs/security-training-image-freiheit93-rootfs.spdx.json | wc -l
23938
\end{bashinput}

Analyzing this SBoM manually is not a good option, even though the system we are using is a rather bare-bones one.

\section{Analyzing the SBoM}

In this section, we are going to use \href{https://sbom-cve-check.readthedocs.io/en/latest/introduction.html}{sbom-cve-check}
to analyze the SBoM generated and determine which CVEs affect our target's rootFS.

\begin{bashinput}
$ sbom-cve-check --sbom-type spdx3 -S build/tmp-glibc/deploy/images/freiheit93/security-training-image-freiheit93.rootfs.spdx.json -f summary -o cve_summary_novex
\end{bashinput}

\begin{bashinput}
INHERIT:remove = "create-spdx"
INHERIT:append = " create-spdx-3.0"
INHERIT:append = " vex"
\end{bashinput}

\begin{bashinput}
$ bitbake security-training-image
\end{bashinput}

This should have created the VEX manifest in \code{build/tmp-glibc/deploy/images/freiheit93/security-training-image-freiheit93.rootfs.json},
which we can now use to add context to \code{sbom-cve-check}:

\begin{bashinput}
$ sbom-cve-check --sbom-type spdx3 \
	-S build/tmp-glibc/deploy/images/freiheit93/security-training-image-freiheit93.rootfs.spdx.json \
	--yocto-vex-manifest build/tmp-glibc/deploy/images/freiheit93/security-training-image-freiheit93.rootfs.json \
	-f summary -o cve_summary_vex
\end{bashinput}

Compare both outputs to find out the impact of the annotations built into Yocto.

% CVE-2024-58251

Take a look at \href{https://nvd.nist.gov/vuln/detail/cve-2024-9287}{CVE-2024-9287}. Does it actually apply here?
If not, try to correct \code{sbom-cve-check}'s assessment.

The \code{summary} format does not indicate \textbf{why} \code{sbom-cve-check} concluded that this CVE was applicable.
You might want to generate a \code{yocto-cve-check-manifest} output.

\section{Patching a CVE reported by sbom-cve-check}

In this section, we are going to fix one of the CVEs. Specifically, \href{https://nvd.nist.gov/vuln/detail/CVE-2026-33317}{CVE-2026-33317}.

Take a look at the CVE, patch it, and make sure it no longer appears in \code{sbom-cve-check}'s report.
